\documentclass{article}
\usepackage[utf8]{inputenc}
\usepackage{geometry}
\usepackage{xcolor}
\geometry{a4paper, margin=1in}

\title{Evaluación de Respuestas: José}
\author{evaluado por Prof. CESAR RODRIGUEZ}
\date{\today}

\begin{document}

\maketitle

\section*{Análisis del Cuestionario}

A continuación se presenta la evaluación de tus respuestas a la "Evaluación de la Semana 1: 
Consultas DNS (TXT/SOA)".

\subsection*{Pregunta 1: Concepto de Registro SOA}
\begin{itemize}
    \item \textbf{Tu Respuesta (transcrita):} \textit{"Este tipo de Registro DNS nos 
    puede proporcionar información del administrado como el correo y otras información de autoridad 
    como el nombre que esta recibe que es inicio de autoridad, tambien se consiguen resgistro 
    autoritarios."}
    \item \textbf{Evaluación:} \textbf{\textcolor{orange}{Regular}}.
    \item \textbf{Comentarios:} Aquí dentificas correctamente que el registro SOA contiene el 
    correo del administrador y está relacionado con la "autoridad" de la zona. Sin embargo, tu 
    respuesta es vaga y omite los componentes más importantes y técnicos de un registro SOA, como 
    el servidor de nombres primario (MNAME), el número de serie (SERIAL) y los temporizadores 
    (REFRESH, RETRY, EXPIRE), que son cruciales para la replicación y coherencia de la zona DNS.
\end{itemize}

\subsection*{Pregunta 2: Registros SPF y Seguridad}
\begin{itemize}
    \item \textbf{Tu Respuesta (transcrita):} \textit{"El Presgistro Txt, busca los registro de 
    texto como Correos y otros datos almacenado en txt, el Spf, Son los almacenamientos de Correos 
    que son usados por nombre de dominios."}
    \item \textbf{Evaluación:} \textbf{\textcolor{red}{Errada}}.
    \item \textbf{Comentarios:} Tu respuesta es muy confusa y no explica qué es un registro SPF 
    ni qué problema de seguridad mitiga. No mencionas que SPF (Sender Policy Framework) es un 
    mecanismo para prevenir la \textbf{suplantación de correo electrónico (email spoofing)}. 
    Tu descripción de "almacenamientos de Correos" es incorrecta; un registro SPF no almacena 
    correos, sino que define qué servidores de correo están autorizados para enviar emails en 
    nombre de un dominio.
\end{itemize}

\subsection*{Pregunta 3: Excepción en \texttt{dnspython}}
\begin{itemize}
    \item \textbf{Tu Respuesta (transcrita):} \textit{"Al no encontrar información del Registro DNS, 
    entrato este envia por Pantalla que no fue encontrado un Registro DNS para el Dominio que esta 
    consultando."}
    \item \textbf{Evaluación:} \textbf{\textcolor{red}{Errada}}.
    \item \textbf{Comentarios:} No respondes a la pregunta. La pregunta solicita el 
    \textbf{nombre de la excepción específica} que lanza la biblioteca \texttt{dnspython}. En su 
    lugar, describes el comportamiento de tu propio script (imprimir un mensaje en pantalla). 
    La respuesta correcta era \texttt{dns.resolver.NoAnswer}, que indica que la consulta se 
    resolvió, pero no se encontraron registros del tipo solicitado.
\end{itemize}

\section*{Conclusión de la Evaluación}

Tu desempeño es bajo. Aunque tienes una idea superficial sobre el registro SOA, 
tus conocimientos sobre SPF y el manejo de excepciones en la biblioteca \texttt{dnspython} 
son deficientes o nulos. Las respuestas a las preguntas 2 y 3 demuestran una falta de comprensión 
tanto del problema de seguridad como de la herramienta de software utilizada.

\begin{itemize}
    \item \textbf{Pregunta 1:} Regular
    \item \textbf{Pregunta 2:} Errada
    \item \textbf{Pregunta 3:} Errada
\end{itemize}

\textbf{CALIFICACIóN: 3.33 en base a 10 punto.}
\end{document}